\par\medskip\noindent\begin{minipage}{\linewidth}\centering
\begin{tikzpicture}\begin{axis}[width=.88\linewidth,height=6.4cm,axis lines=middle,grid=major,grid style={gray!25},tick label style={font=\normalsize},label style={font=\normalsize},clip=false,axis equal image,xmin=-.5,xmax=5,ymin=-.5,ymax=5.5,xlabel={$x$},ylabel={$y$},xtick={0,1,2,3,4,5},ytick={0,1,2,3,4,5}]\addplot[mxink,thick,domain=0:2.65]{2*x};\addplot[mxmint,very thick] coordinates {(4,3)(2,4)};\draw[dashed,gray](axis cs:4,0)--(axis cs:4,3)--(axis cs:0,3);\addplot[only marks,mark=*] coordinates {(4,3)(2,4)};\draw(axis cs:2.1342,4.2683)--(axis cs:2.4025,4.1342)--(axis cs:2.2683,3.8658);\node[above right,fill=white,inner sep=2pt]at(axis cs:4,3){$P(4;3)$};\node[left,fill=white,inner sep=2pt]at(axis cs:2,4){$H(2;4)$};\node[above right,fill=white,inner sep=2pt]at(axis cs:2.4,4.8){$d:y=2x$};\end{axis}\end{tikzpicture}
\par\medskip Le segment PH est perpendiculaire à d. Les deux axes ont la même échelle ; les coordonnées de P et H se lisent directement.\end{minipage}\par\medskip
